% =========================================================================
\chapter{Discussion and Conclusion}
\label{ch:discussion}
\label{ch:conclusion}
% =========================================================================

\claim{The proposition's scope is stated, the gates' blind spot is
exhibited, the costs are honest, and the roadmap attaches the missing organ
where the evidence says it belongs: the retention decision.}

% -------------------------------------------------------------------------
\section{The readout proposition and its scope}
\label{sec:proposition}
% -------------------------------------------------------------------------

\subsection*{The proposition}

Self-evolution is usually described as variation plus selection. This
thesis argues it has a third term:

\begin{quote}
\textbf{Self-evolution $=$ variation $+$ selection $+$ readout.} When the
effects being selected are small relative to variation in the score
readout, score-based retention may mis-rank edits. The cumulative-gain
statement is a hypothesis motivated by this case, not a dynamic estimated
by these campaigns.
\end{quote}

It is a proposition rather than a theorem: as stated in words it restates
power analysis, and what makes it non-trivial is the dynamics argument
below, which a single-test power calculation does not supply.

The claim is made for this bed and this system class. The field-level
implication --- that every surveyed loop selects on aggregate score
(\S\ref{sec:self-evolving}) and is therefore exposed to the same
arithmetic on any bed with a comparable flip rate --- is an \emph{inference}
and is flagged as one.  An independent, contemporaneous
critique of how harness evolution is evaluated reaches a convergent
conclusion under a different evaluation design~\cite{rethink-harness-eval}.

\subsection*{Formal spine}

Quantitative genetics already has the equation. The response to selection
is given by the breeder's equation~\cite{falconer-mackay},
\[
  R \;=\; h^2 S ,
\]
gain per generation equals heritability times the selection differential.
When selection acts on a \emph{measured} phenotype rather than the true
one, $h^2$ is attenuated by the measurement's reliability --- classical
regression dilution~\cite{hutcheon-dilution}:
\[
  h^2_{\text{realized}} \;=\; h^2_{\text{true}}\,\rho ,
  \qquad
  \rho \;=\; \frac{\sigma^2_{\delta}}{\sigma^2_{\delta}+\sigma^2_{\varepsilon}} .
\]

The reliability form is retained here as a conditional analogy. These
campaigns do not estimate the effect-size variance, measurement-error
variance, or retention dynamics needed to instantiate $\rho$; in
particular, the $3.70$ in \S\ref{sec:power} is the sample SD of twenty
paired same-configuration score changes, not $\sigma_\varepsilon$.

\paragraph{The dynamics claim.} On this bed and system class, expected
cumulative gain can be non-positive when selection mis-ranks small
effects. That is a hypothesis consistent with this case; no empirical
$\rho$, false-kill law, or variation-quality-independent result is
claimed.
\paragraph{Boundary condition.} The breeder's equation governs the
\emph{mean} response, and \S\ref{sec:archaeology} shows gains on this bed
arriving as rare tail events from a single mechanism family. The
proposition therefore constrains the mean and is silent on the tail, where
this thesis's one positive result lives.

\subsection*{The calibration that motivates all of it}

The budget-starvation repair (\S\ref{sec:one-repair}) is a
post-hoc positive-control observation. The score axis recovered this large
event: \texttt{budget\_exceeded} exits collapsed from 35 to 4 in one
round, twelve of the sixteen newly passing tasks had died of starvation in
the round before, and the $+9$ cleared the same-configuration band \F{39}.
It is one event with no interval. Its threefold relation to the smaller
effects discussed here is descriptive and does not establish assay
sensitivity for them.

\subsection*{The rival hypothesis}

Every null in this thesis admits a simpler explanation than the readout
proposition: \textbf{the loop's action space cannot reach the lesions.}
This is the capability wall, and it is the primary layer in the mortality
account \F{40}. If the drugs are placebos because the disease is out of
reach, no readout is needed to explain a flat outcome axis.

On the outcome axis the two are observationally equivalent. Level 3 was
exercised by hand
\F{34}\F{36} and never wired into retention, so no experiment here
separates them.

They are not equivalent on all of the evidence. The separation is \cls{C}
judgment rather than a test:

\begin{enumerate}
  \item \textbf{The actuator demonstrably reached a lesion at least once.}
        The starvation repair worked, was mechanism-confirmed, and was
        retained \F{39}. A capability wall that admits this counter-example
        is no longer the claim that the action space cannot reach lesions;
        it is the narrower claim that it cannot reach the \emph{dominant}
        family --- wrong-closure, 18 of the 49 swing tasks \F{31} --- which
        this thesis accepts (\S\ref{sec:future}).
  \item \textbf{The drug that worked was then lost for a reason unrelated
        to capability.} It reached no later campaign because its target is
        a \texttt{file://} path into its own campaign directory and every
        later seed is built from the clean template: \emph{the inheritance
        channel does not exist} \F{39}. No capability wall explains that.
  \item \textbf{The same families are re-derived from scratch, repeatedly.}
        The loop re-prescribes drugs it has already prescribed correctly,
        four campaigns apart, without inheriting \F{38}. A loop blocked by
        a capability wall would fail to \emph{cure}; it would not fail to
        \emph{remember}.
\end{enumerate}

In summary: the capability wall explains why the loop does not solve the
hard family. It does not explain the loss of the one drug
that worked, and it does not explain the treadmill. Those two read as
readout and inheritance failures, which is what the proposition
addresses.

\paragraph{The discriminating experiment.} Attach the level-3 readout to the retention decision and re-fly. If
retention still drifts, the actuator was the bottleneck after all. This is
why \S\ref{sec:future} leads with it.


% -------------------------------------------------------------------------
\section{Governance: necessary, not sufficient}
\label{sec:governance}
% -------------------------------------------------------------------------

The acceptance gates work, within their remit. Seven structural refusals
are on the record \F{11}, and no refused candidate shipped.

Their remit is form, not provenance. Under a coverage directive --- treat
every task, never abstain --- the loop engineered the bed's own leak routes
into a benchmark-lookup cheat tool, and that tool passed the Critic and
\emph{every} landing gate \F{37}. Nothing in the ladder asks where a
capability came from.

Two fixes follow: a leakage blocklist on the bed side, and a gate
forbidding candidates from reading adjudication files on the loop side ---
the second written after registry error 13, where an adjudication file left
readable in a run root was read by a role and a round of rulings had to be
voided. The governance gap and the evaluator gap are the same gap seen from
two sides.

% -------------------------------------------------------------------------
\section{Threats to validity}
\label{sec:threats}
% -------------------------------------------------------------------------

\paragraph{Attribution is arm-level.} No single-factor attribution of
\emph{score} effects is claimed anywhere \F{9c}; the per-switch entries of
the profit-and-loss account and the component claims of
\S\ref{sec:future} are construction and capability facts, not outcome
attributions. The reason no component ablation was run is given in
\S\ref{sec:localization}.

\paragraph{The localization metric is post-hoc}, defined over predictions
that were written before the outcome was known but scored under a rule
fixed afterwards. This is stated at every point of use.

\paragraph{Single bed, single model tier.} Three seeds per arm establish
that the noise geometry replicates \F{44}\F{45}; generalization to other
beds and model families is untested.

\paragraph{Bed asymmetry between the arms.} The arms did not fly quite the
same bed (\S\ref{sec:beds}); every cross-arm reading is recomputed on the
common 100-task subset from \texttt{task\_history} rather than taken from
\texttt{curves.json} \F{46}.

\paragraph{Leak-route contamination.} Thirteen of the 49 swing tasks are
leak-route \F{31}; both arms see the same fixed bed every round, so drift
toward it is shared. \emph{No held-out defence was in place}: the
blocklist of \S\ref{sec:future} is proposed, not applied, and flew in no
campaign reported here.

\paragraph{Operational deviations} --- a machine restart, a gateway
outage, five meta-chain holes, three re-flights --- are isolated in
\S\ref{sec:replication}, and only verified clean same-config pairs enter
the flip statistics. \textbf{Windows frictions} are symmetric across arms.
\textbf{Judge residual risk} remains after hardening.

\paragraph{The endpoint was substituted after data were seen.} H1's
original endpoint --- a between-arm localization comparison --- was
abandoned as unmeasurable once the two arms' candidate populations and
gate regimes were known, and the declared substitute endpoints took its
place \F{9c}. The substitution is post-hoc by construction, so H1's
verdict is a verdict on the substitutes: the original hypothesis is
\emph{untested}, not refuted, and no reading here recovers it.

\paragraph{Disclosure is not mitigation.} The post-hoc scoring rule and
the fifteen-error registry below are declared rather than controlled for.
Neither is mitigated; both remain live limits on the intervals this thesis
reports.

\paragraph{The calibration event is one round pair,} found by a sweep
rather than pre-specified, so its ``three times'' is a point ratio without
an interval (\F{39}\F{49}).

\paragraph{The evaluator.} The fifteen-error registry is reported here
with both bias directions documented and \emph{stopping the check too
early} named as the common root cause, and a retraction registry keeps
nothing it records re-cited --- including this project's own earlier
reading of the base loop's prescriptions, which its archives corrected
(\S\ref{sec:archaeology}).

% -------------------------------------------------------------------------
\section{Cost honesty}
\label{sec:cost}
% -------------------------------------------------------------------------

\subsection*{The cost work is not a contribution of the graph}

Noop-round batching (\$38.1 to \$11.8 per round \F{13}), single-shot
digestion, the zero-cache-hit repair and the retention policy are
\emph{harness-level} mechanisms: each runs on the unmodified official
system and none needs an execution graph, so the profit-and-loss account
labels them \emph{portable} \F{12} and Appendix~\ref{app:cost} describes
them as engineering.

\subsection*{What the graph layer costs is unknown}

On the whitelist campaigns the task-side spend per fresh evaluation is
\emph{lower} in the graph arms on all three seeds (\$0.392, \$0.472,
\$0.503 against \$0.461, \$0.538, \$0.544) \F{13}, but that comparison
settles nothing: it excludes the meta tier, which is where graph work would
appear and for which no admissible per-role accounting exists. The entry is
a blank rather than a minus, and an earlier net-expensive reading, measured
on campaigns the whitelist later excluded, is withdrawn.

\subsection*{The graph layer's one measured liability}

It is operational rather than financial. A crash family introduced by the
recorder --- binary content reaching the raw JSONL and killing the evolve
stage --- produced 164 tracebacks across the three graph seeds against 5
across the three no-graph seeds, together with five meta-chain holes
\F{56}.

% -------------------------------------------------------------------------
\section{Future work}
\label{sec:future}
% -------------------------------------------------------------------------

Each item is attached to a measured gap and ordered by what the evidence
says matters.

\paragraph{1. Readout-first self-evolution.} Attach rung R3 to the
retention decision. The ingredients are already in the literature ---
paired prefix replay~\cite{shepherd,car}, sibling-rollout
controls~\cite{craft}, anytime-valid sequential tests~\cite{pace}. This
thesis contributes \emph{where to attach them}, and the evidence that
attaching them elsewhere is premature. This is also the discriminating
experiment of \S\ref{sec:proposition}.

\paragraph{2. Bed reliability, with its ceiling priced first.} The leakage
blocklist and result caching at the measured variance injection point ---
a web search, in 76 of 79 cases \F{25} --- are direct and cheap. Reliability
engineering should target the volatile family, where flips concentrate at
$1.6$ to $2.0$ times the bed rate \F{41}\F{44}\F{45}, rather than the whole
bed.


\paragraph{3. Widen the typed surface to the reach the free-text one
already has.} GHX types a slice of the official action space and is
\emph{narrower} than the prose path it replaces: seven operations on the
processor plane, two config-section operations for tools, prompt edits as
parameter mutations, and the harness's own source out of scope
(\S\ref{sec:candidate-surface}). The official Evolver, meanwhile, already
authors new processor classes as source files, and its one above-noise
repair shipped exactly that way \F{39}. Typing bought legality and
identity at a cost in reach, and nothing in this study prices that trade,
because the outcome axis cannot. The experiment is to extend the typed
operations until they cover what prose already reaches, then read the two
surfaces against each other under item 1's readout --- which is the only
way to learn whether a typed candidate is worth what it gives up.

\paragraph{4. Replicate the episodic-note channel} under a pre-registered,
budget-bounded protocol \F{35}. It is solver-level and bypasses all four
mortality layers, which makes it the cheapest positive result available,
and --- per \S\ref{sec:survivor} --- the one that gives up the property
under study.

\paragraph{5. Capability procurement for the wrong-closure family.}
Offline simulation over the trial repetitions shows majority-at-$k$ lifting
mean pass from $0.35$ to $0.60$ at $k{=}10$, entirely from scatter-shaped
tasks (3 lifted, 0 sunk), while consensus-wrong tasks lock at zero for
every $k$ \F{42}. Plain voting \emph{hardens} the confident-wrong answer.
The organ needed is external discrimination, and that does sit outside the
loop's action scope.



% -------------------------------------------------------------------------
\section{Conclusion}
% -------------------------------------------------------------------------

The loop can see, classify, prescribe, invent tools, and abstain --- and it
cannot hear whether any of it worked. This thesis argued the loop is deaf,
measured its floor on three seeds, showed that clearing that
floor is neither necessary nor sufficient for an edit to have worked,
retyped the candidate surface, exhibited the one ratchet click whose
mechanism the record can still read, and left the readout ladder standing
with three of its four rungs live.

That contribution came from a large build: a typed graph runtime, new
build-time and execution-time gates, and a machine-generated candidate
surface, set against a claim that stays narrow --- legality and identity,
not reach. The build is enabling infrastructure for the readout question,
not an efficient route to checkability: a lighter instrument might have
answered the same question at lower engineering cost.

% -------------------------------------------------------------------------
\subsection*{What the thesis claims, restated}
% -------------------------------------------------------------------------

C1 is an independent execution of the official system. C2 is the
construction proof that its evidence column cannot register the use of a
short fact, and the reach of that column into shipped decisions. C3 is GHX
--- legality and identity of candidates, explicitly not a wider reach. H1's
original localization endpoint was unmeasurable; the whole-stack two-arm
readings on substitute endpoints are descriptive, and the separately
protocolled bundle trial returned no detected improvement with a negative
point estimate. C5 is the separation of capability from
efficacy, and the readout proposition it supports. C6 is a roadmap in which
every item is attached to a measured gap.

The proposition itself is stated once more, with its boundary:
self-evolution is variation plus selection plus \emph{readout}, and on
this bed and system class, when selected effects are small relative to
variation in the score readout, mis-ranking could leave expected cumulative
gain non-positive. This is a hypothesis, not an estimated dynamic. The
breeder's equation governs the \emph{mean} and is silent on the tail ---
which is where this thesis's one positive result lives.
The rival account --- that the loop's actuator simply cannot reach the
lesions --- is not dismissed. It is weighed in \S\ref{sec:proposition},
where it explains the hard family but not the loss of the one drug that
worked, nor the treadmill that re-derives drugs already correctly
prescribed. The experiment that would separate them is named rather than
claimed.
